\Idx{阿贝尔簇}{428}
\Idx{伴随表示}{70, 90}
\Idx{容许合成律}{129}
\Idx{仿射群}{20, 403}
\Idx{\(\mathbf{Alf}_{/k}\) (有限长 \(k\)-代数范畴)}{515}
\Idx{Hopf 代数}{21, 427, 546, 565, 582}
\Idx{Lie 代数}{}
\Idx{\quad 形式群的}{559}
\Idx{\quad 群概形的}{89, 90}
\Idx{中心化子 (正规化子) 的 Lie 代数}{}
\Idx{\quad 子群的}{94}
\Idx{\quad 子模的}{97}
\Idx{截断形式幂级数代数}{589}
\Idx{\(G\) 的分布代数}{449}
\Idx{受限包络代数}{473, 476, 575}
\Idx{\(G\) 的无穷小代数}{449}
\Idx{\(\mathbf{Alp}_{/k}\) (预有限 \(k\)-代数范畴)}{514}
\Idx{\(\alpha_{p,k}\)}{99}
\Idx{反仿射 \(k\)-群}{429}
\Idx{\(a(P)\) (预层 \(P\) 的伴随层)}{209}
\Idx{\(\operatorname{Aut}(X)\)}{9, 12}
\Idx{余交换双代数}{544}
\Idx{双代数 (bigebra; 另见 bialgebra)}{20}
\Idx{良好 \(\mathcal O_S\)-模}{81}
\Idx{良好群函子}{83}
\Idx{局部 Boolean 空间}{363}
\Idx{Cantor 空间}{364}
\Idx{特征子群}{377}
\Idx{交换 \(k\)-群范畴}{}
\Idx{\quad 仿射}{555}
\Idx{\quad 代数}{323}
\Idx{\quad 仿射代数}{323}
\Idx{\quad 拟紧}{327}
\Idx{\(\widehat{\mathcal C}\)}{1}
\Idx{\(\widehat{\mathcal C}_{/S}\)}{3}
\Idx{\(\widehat{\mathcal C}^{\,d}_{/S}\), \(\widehat{\mathcal C}^{\,b}_{hS}\)}{3}
\Idx{中心子群}{17}
\Idx{中心化子}{15, 366}
\Idx{\(\underline{\operatorname{Centr}}_G(X)\)}{15}
\Idx{\(\underline{\operatorname{Centr}}(G)\)}{17}
\Idx{Chevalley 定理}{420, 422, 428}
\Idx{形式簇 \(X\) 的余代数 \(H(X)\)}{530}
\Idx{余交换余代数}{455}
\Idx{群余代数}{457, 477}
\Idx{余代数 (cogebra; 另见 coalgebra)}{409}
\Idx{Cohen--Macaulay}{294}
\Idx{换位子}{}
\Idx{\quad 换位子层}{388}
\Idx{\quad 换位子群}{386}
\Idx{余模}{28, 409}
\Idx{单位连通分支 \(G^0\)}{299, 308, 345}
\Idx{局部分量 \(A_m\), \(M_m\)}{502, 511}
\Idx{条件 (E)}{59}
\Idx{几何连通}{297}
\Idx{余法层}{46, 147}
\Idx{一对箭头的余核}{178, 186, 250, 310, 492, 519}
\Idx{常对象}{10}
\Idx{\quad \((\mathrm{Sch})\) 中的}{20}
\Idx{等价对}{254}
\Idx{覆盖态射或覆盖族}{200}
\Idx{协变双代数或协变代数}{545, 547}
\Idx{筛}{196}
\Idx{\(\operatorname{Lie}(G/S)\) 上的括号}{85}
\Idx{\(c(X,Y,f)\)}{152}
\Idx{\(\Delta^{(1)}_{X/S}\)}{105}
\Idx{\(S\)-态射 \(Y\to X\) 的 \(S\)-导子}{443}
\Idx{不变导子}{89}

\Idx{下降}{}
\Idx{\quad 下降数据}{182}
\Idx{\quad 下降态射}{183}
\Idx{投射性的下降}{413}
\Idx{有效下降态射}{183}
\Idx{阶数 \(\leq n\) 的 \(S\)-偏差}{441}
\Idx{纤维维数}{350}
\Idx{\(D_{\mathcal O_S}(M)\)}{52}
\Idx{Cartier 对偶}{461, 547}
\Idx{\(d(Y_0,Y)\)}{148}
\Idx{仿射包络 \(G_{\mathrm{af}}\)}{405, 427}
\Idx{满态射}{}
\Idx{\quad 有效}{178}
\Idx{\quad 普遍有效}{178}
\Idx{\quad 普遍}{177}
\Idx{等变对象与等变模}{38}
\Idx{商赋环空间}{250}
\Idx{\(u\) 处的切空间}{56}
\Idx{形式群的齐性空间}{582, 590}
\Idx{本质自由}{367}
\Idx{平展}{}
\Idx{\quad 形式平展代数}{539}
\Idx{\quad 形式群}{556}
\Idx{\quad 形式簇}{539}
\Idx{零延拓}{416}
\Idx{预层的伴随层}{208}
\Idx{相对层}{222}
\Idx{商层}{217, 259, 272, 277, 280, 283, 286, 287, 393, 492, 493, 537, 552}
\Idx{向量丛 \(V(F)\)}{25}
\Idx{主齐性丛}{231}
\Idx{切丛}{55}
\Idx{忠实作用}{17}
\Idx{函子 \(\operatorname{Aut}(X)\)}{9, 12}
\Idx{\((\mathrm{Sch})\) 上的环函子 \(\mathcal O\)}{22}
\Idx{函子 \(\operatorname{Hom}(X,Y)\)}{7}
\Idx{函子 \(\operatorname{Hom}_{Z/S}(X,Y)\)}{50}
\Idx{函子 \(\prod_{Z/S}Y\)}{51, 368, 373}
\Idx{函子 \(X\mapsto X_e\)}{542}
\Idx{模函子 \(V,W\)}{24}
\Idx{形式概形}{503, 517}
\Idx{形式簇}{518}
\Idx{形式平展代数}{539}
\Idx{形式齐性空间}{246}
\Idx{形式主齐性}{102, 230}
\Idx{Frobenius 态射}{}
\Idx{\quad 绝对, \(\operatorname{fr}(S)\)}{461, 574}
\Idx{\quad 相对, \(\operatorname{Fr}(X/S)\)}{461, 573}
\Idx{\(\Gamma^*(M)\)}{522}
\Idx{\(\mathbb G_a\)}{21}
\Idx{\(\mathbb G_m\)}{22}
\Idx{\(G\)-\(\mathcal O\)-模}{19}
\Idx{\(G\)-\(\mathcal O_S\)-模}{27, 43}
\Idx{完美域上的 \(G_{\mathrm{red}}\)}{291}
\Idx{不完美域上的 \(G_{\mathrm{red}}\)}{297}
\Idx{范畴中的群}{12, 13}
\Idx{带算子群的对象}{15}
\Idx{平展群 \(G/G^0\)}{323}
\Idx{域上的平滑群}{296}
\Idx{纤维连通的群}{353}
\Idx{Dedekind 概形 \(S\) 上的仿射群}{430, 431}
\Idx{可对角化群}{23}
\Idx{\quad 其上同调}{37}
\Idx{域 \(k\) 上的拟紧群}{326}
\Idx{群胚}{251}
\Idx{\(H\)-集合}{62}
\Idx{\(h_A\), 函子 \(\operatorname{Hom}_{\mathbf{Alp}/k}(A,-)\)}{515}
\Idx{高度 \(\leq 1\), 对 \(S\)-群及形式群而言}{489, 492, 579}
\Idx{高度 \(\leq n\), 对 \(S\)-群或形式群而言}{464, 575, 590}
\Idx{\(H^\vee_0(S,P)\)}{205}
\Idx{\(h^c_M\), 函子 \(\operatorname{Hom}_c(M,-)\)}{506}
\Idx{齐性空间}{246}
\Idx{交叉同态}{74}
\Idx{\(\operatorname{Hom}(X,Y)\)}{7}
\Idx{\(\operatorname{Hom}_{\mathcal O}(F,F')\)}{19}
\Idx{\(\operatorname{Hom}_{Z/S}(X,Y)\)}{50}
\Idx{\(\operatorname{Hom}_{(Z/S)\text{-gr.}}(X,Y)\)}{74}
\Idx{\(H(X)\)}{530}
\Idx{\(h_X\)}{1}
\Idx{正则浸入}{149, 161, 597}
\Idx{无穷小 \(S\)-群或形式群}{484, 557, 569}
\Idx{无穷小}{}
\Idx{\quad 自同构}{98}
\Idx{\quad 自同态}{68, 69}
\Idx{完全交}{597}
\Idx{局部完全交}{149, 160}
\Idx{不变子群}{17}
\Idx{不变元子对象}{15, 28}
\Idx{\(I_S(M)\)}{52}
\Idx{Jacobson 公式}{471, 472}
\Idx{\(\operatorname{Ker}(f)\)}{17}
\Idx{Koszul 复形}{149}
\Idx{\(\mathbf{LF}(A)\), 有限长 \(A\)-模范畴}{506}
\Idx{自由作用的群}{188}
\Idx{\(\operatorname{Lie}(X/S)\)}{63, 70}
\Idx{\(\operatorname{Lie}(X/S,M)\)}{63, 70}
\Idx{\(\operatorname{Lie}(f)\), \(f\) 的导出态射}{70}
\Idx{\(\operatorname{Lie}_0(X/S,M)\)}{77}
\Idx{\(\operatorname{Lie}(G/S)\)}{89, 90}

\Idx{\(\operatorname{Lie}(\mathbb G_{m,S}/S)\), \(\operatorname{Lie}(D_S(M)/S)\)}{91}
\Idx{\(\operatorname{Lie}({}_nG)\), \(g\mapsto g^n\) 的导出态射}{64}
\Idx{群概形的逆极限}{396}
\Idx{仿射代数群的线性性}{}
\Idx{\quad 在维数 \(\leq2\) 的正则 \(S\) 上平坦}{433}
\Idx{\quad 在域上}{417}
\Idx{\(G/\operatorname{Fr}^nG\) 的平滑性}{493}
\Idx{特征零的域 \(k\) 上 \(G\) 的平滑性}{336, 571}
\Idx{\(\mathrm{LP}\)}{205}
\Idx{\(L^u_{X/S}\)}{56}
\Idx{\(L^u_{X/S}(M)\)}{56}
\Idx{\(L_X\)}{104}
\Idx{\(L^0_X\), \(L'_X\)}{114}
\Idx{\((M)\)-有效性}{193}
\Idx{\(M^\dagger,N^*\)}{507}
\Idx{单态射与闭浸入}{301, 327, 335}
\Idx{态射 \(g\mapsto g^n\) 平展}{495}
\Idx{态射 \(g\mapsto g^n\) 为零}{496}
\Idx{\(\mu_n\), \(n\) 次单位根}{23}
\Idx{Noether 同构定理}{235}
\Idx{\(S\) 上的对偶数概形}{52}
\Idx{正规化子}{15, 366}
\Idx{范数}{}
\Idx{\quad 可逆层的}{267}
\Idx{\quad 有限局部自由 \(A\)-代数的}{261}
\Idx{\(\underline{\operatorname{Norm}}_G(X)\)}{15}
\Idx{核}{17}
\Idx{\(N_{Y/X}\), \(Y\) 在 \(X\) 中的余法层}{46, 147}
\Idx{\(\Omega^1_{X/S}\)}{54, 89, 105}
\Idx{\(\omega^1_{X/S}\)}{89}
\Idx{\(\mathcal O\)-模}{18}
\Idx{平坦 \(\mathcal O_k\)-模}{520}
\Idx{微分算子}{441}
\Idx{\(G\) 上的不变微分算子}{449}
\Idx{\(p\)-Lie 代数}{470}
\Idx{\quad \(S\)-群的}{480}
\Idx{\quad 形式群的}{575}
\Idx{\(\mathbf{PC}(A)\), 伪紧 \(A\)-模范畴}{503}
\Idx{\(\prod_{T/S}F\) 的上同调}{124}
\Idx{\(\prod_{Z/S}Y\)}{51}
\Idx{点化余代数或点化双代数 (pointed)}{556, 565}
\Idx{几何逐点不可约}{306}
\Idx{预层范畴}{1}
\Idx{分离预层}{202}
\Idx{预等价关系}{249, 253}
\Idx{预拓扑}{201}
\Idx{本原元素}{459, 561, 566}
\Idx{完备化张量积 \(\widehat{\otimes}\)}{508, 509, 511, 516, 519}
\Idx{预有限 \(k\)-代数}{514}
\Idx{伪基}{507}
\Idx{伪紧 (pseudocompact)}{}
\Idx{\quad 环}{501}
\Idx{\quad 模}{503}
\Idx{可平方态射}{179}
\Idx{拟截面}{268}
\Idx{拟分离概形}{307}
\Idx{\(S\) 上的拟分离概形}{363, 403, 409}
\Idx{对群胚取商}{}
\Idx{\quad 有限平坦}{259}
\Idx{\quad 平坦, 不必固有}{277}
\Idx{\quad 固有平坦}{272}
\Idx{Artin 局部环 \(A\) 上的商 \(G/H\)}{311, 315}
\Idx{\(\mathbf{Vaf}_{/k}\) 中的商}{534, 537}
\Idx{对群概形取商}{282}
\Idx{\(r(A)\), \(A\) 的 Jacobson 根}{502}
\Idx{加细}{199}
\Idx{等价关系}{186}
\Idx{\quad 有效}{191}
\Idx{\quad 普遍有效}{191}
\Idx{可表性}{}
\Idx{\quad \(\prod_{Z/S}Y\) 的}{368, 373}
\Idx{\quad 中心化子的}{371, 376}
\Idx{\quad 正规化子的}{371, 376}
\Idx{可表函子}{2}
\Idx{可解群或幂零群}{390}
\Idx{Weil 标量限制}{51}
\Idx{内拟紧}{307}
\Idx{\(S\)-\(H\)-函子}{62}
\Idx{\((\mathrm{Sch})\), \((\mathrm{Sch})/S\), \((\mathrm{Sch}_{/S})\)}{20}
\Idx{概形稠密}{374}
\Idx{概形占优}{325, 326, 378, 427}
\Idx{半直积}{16, 17, 552}
\Idx{分离 (域上的每个群都如此)}{292}
\Idx{\(\widehat S_k(E)\), 完备化对称代数}{523}
\Idx{由 \(f:X\to G\) 生成的子群}{378}
\Idx{\(\operatorname{Spec}^*(U)\)}{456}
\Idx{\(\operatorname{Spf}^*(C)\)}{530}
\Idx{\(\underline{\operatorname{Stab}}_G(x)\)}{15}
\Idx{稳定子}{15, 44}
\Idx{严格有理点}{293}
\Idx{拓扑}{}
\Idx{\quad \((\mathrm{\acute etfg})\)}{244}
\Idx{\quad \((\mathrm{fpqc})\), \((\mathrm{fppf})\), \((\mathrm{\acute et})\), \((\mathrm{\acute etf})\)}{241}
\Idx{\quad 典范}{204}
\Idx{\quad 混沌 (粗)}{211, 245}
\Idx{\quad Zariski}{236}
\Idx{\quad 平坦}{537, 552}

\Idx{拓扑}{199}
\Idx{拓扑自由}{507}
\Idx{拓扑幂零}{502}
\Idx{拓扑平坦态射}{527}
\Idx{拓扑平坦形式簇}{529}
\Idx{挠子}{231}
\Idx{传送子}{366}
\Idx{严格传送子}{366}
\Idx{甚好群函子}{88, 89, 459}
\Idx{\(T_{X/S}\)}{55}
\Idx{\(T_{X/S}(M)\)}{55}
\Idx{\(T'_{X/S}(M)\)}{77}
\Idx{\(k\) 上的乘法型群}{558}
\Idx{幺幂 \(k\)-群}{558, 563}
\Idx{\(\Upsilon(A)\), \(A\) 的极大开理想}{502}
\Idx{\(\mathbf{Vaf}_{/k}\), \(k\) 上的形式簇范畴}{518}
\Idx{\(\mathbf{Vaf}_{/S}^{f\flat}\)}{525}
\Idx{\(\mathbf{Vaf}_{/k}^{\acute et}\)}{539}
\Idx{阿贝尔簇}{428}
\Idx{余代数 \(C\) 的形式簇 \(\operatorname{Spf}^*(C)\)}{530}
\Idx{Verschiebung}{468}
\Idx{\(V^f_k(N)\)}{520}
\Idx{\(V(\Omega^1_{X/S})\)}{55}
\Idx{\(V^f_k(E)\)}{524}
\Idx{\(V^{f,0}_k(E)\)}{524}
\Idx{\(X\), 概形 \(X\) 的底集}{249}
\Idx{\(X^+=\prod_{S_J/S}X_J\)}{102}
\Idx{\(X_e\), 与 \(X\) 伴随的平展形式簇}{542}
\Idx{\(\widehat X_{/S}\), 与 \(S\)-概形 \(X\) 伴随的形式簇}{525}
\Idx{Yoneda 引理}{1}
